\section{Исследовательская часть работы}
\label{sec:Chapter2} \index{Chapter2}

В данной главе обосновывается архитектура разрабатываемого канала диагностики и
описывается алгоритм кольцевого буфера.

\subsection{Требования к каналу диагностики}
\label{sec:ch2-requirements}

Из постановки задачи и обзора альтернатив следуют следующие требования:
\begin{enumerate}
    \item канал должен работать на ранних стадиях UEFI, когда оперативная память
    может быть ещё недоступна;
    \item после инициализации памяти запись должна выполняться без побайтового
    выхода в эмулятор;
    \item несколько источников в гостевой системе должны иметь возможность писать
    сообщения в общий поток;
    \item пользовательский режим должен иметь путь записи с минимальными
    накладными расходами на обращения к драйверу;
    \item эмулятор должен объединять сообщения гостевой системы со своими отладочными сообщениями в
    хронологическом порядке;
    \item переполнение буфера должно быть явно обнаруживаемым и управляемым.
\end{enumerate}

Эти требования приводят к двухрежимной схеме: синхронная запись через регистр
виртуального устройства используется до появления общей памяти, а основной режим
работы после инициализации памяти строится на кольцевом буфере.

\subsection{Выбор общей памяти и MMIO-уведомления}
\label{sec:ch2-shared-memory}

Один из возможных вариантов архитектуры --- разместить кольцевой буфер в
эмулируемой MMIO-области устройства и позволить гостевой системе писать
сообщения прямо в эту область. Для KVM-подобной виртуализации такой вариант
невыгоден: обращение к эмулируемой MMIO обрабатывается не как обычная запись в
RAM гостя, а как операция устройства. Такая запись может приводить к выходу в
эмулятор и обработке доступа в пользовательском процессе Stratovirt. При
интенсивном логировании это превращает каждую запись в дорогое взаимодействие с
гипервизором.

Поэтому в предлагаемой схеме буфер размещается в RAM гостевой системы. Гость
записывает сообщение обычными инструкциями работы с памятью, а виртуальное
устройство получает только адрес буфера и уведомление о появлении данных. В
результате дорогое обращение к MMIO-регистру требуется не для каждого байта, а
для события завершения записи или принудительного сброса буфера.

Такой выбор также важен для пользовательского режима. Если страницы буфера
отображены в адресное пространство пользовательского процесса, приложение может
формировать запись без дополнительного копирования в промежуточный буфер драйвера.
Драйвер при этом остаётся ответственным за обнаружение устройства, отображение
памяти и выдачу корректного интерфейса доступа.

\subsection{Общая архитектура}
\label{sec:ch2-architecture}

Разработанный механизм включает три основных компонента:
\begin{enumerate}
    \item виртуальное устройство \texttt{LogDev} на стороне Stratovirt;
    \item драйвер и библиотеку логирования для среды UEFI;
    \item драйвер Windows и интерфейсы для пользовательского режима и режима
    ядра.
\end{enumerate}

Каналом взаимодействия между компонентами служит кольцевой буфер, расположенный
в памяти гостевой системы. Виртуальное устройство хранит параметры канала,
предоставляет MMIO-регистры для синхронной записи, адреса буфера и уведомлений,
а также выводит сообщения в лог эмулятора.

Общая архитектура механизма представлена на рисунке~\ref{fig:buffer-architecture}.

\begin{figure}[H]
    \centering
    \includegraphics[scale=1]{buffer.png}
    \caption{Архитектура механизма обмена диагностической информацией}
    \label{fig:buffer-architecture}
\end{figure}

\subsection{Модель виртуального устройства}
\label{sec:ch2-logdev}

Stratovirt --- VMM (Virtual Machine Manager) экосистемы openEuler,
ориентированный на облачные и контейнерные сценарии~\cite{stratovirt_openeuler}.
Как и QEMU, Stratovirt запускает виртуальную машину как отдельный процесс
операционной системы. Внутри процесса выделяются главный поток обработки
событий, потоки vCPU и, при необходимости, отдельные I/O-потоки
(\texttt{iothreads})~\cite{stratovirt_design}.

Устройство \texttt{LogDev} присоединено к шине SysBus и занимает фиксированный
диапазон в карте памяти гостя. Такое решение упрощает доступ к устройству из
UEFI. Устройство добавляется в конфигурацию машины следующим параметром:

\begin{lstlisting}
-device logdevice,id=<logdev_id>[,flush_logs=<true|false>]
[,iothread=<iothread_id>]
\end{lstlisting}

Параметр \texttt{flush\_logs} задаёт поведение при переполнении буфера: ожидать
освобождения места или пропустить сообщение. Параметр \texttt{iothread} позволяет
вынести чтение из буфера в отдельный поток, чтобы снизить нагрузку на главный
цикл эмулятора.

Гостевая ОС взаимодействует с \texttt{LogDev} через MMIO-регистры, перечисленные
в таблице~\ref{tab:logdev-regs}.

\begin{table}[H]
    \centering
    \caption{Регистры MMIO устройства LogDev}
    \label{tab:logdev-regs}
    \begin{tabularx}{\textwidth}{|p{1.2cm}|p{2.0cm}|p{1.6cm}|>{\raggedright\arraybackslash}X|}
        \hline
        Смещ. & Имя & Доступ & Назначение \\
        \hline
        \texttt{0x00} & \makecell[l]{\texttt{SYNC}\\\texttt{\_WR}} & запись &
        Синхронная передача до инициализации буфера \\
        \texttt{0x08} & \texttt{BUFPTR} & \makecell[l]{чтение/\\запись} &
        Физический адрес кольцевого буфера \\
        \texttt{0x10} & \texttt{KICK} & запись &
        Уведомление эмулятора о наличии данных в буфере \\
        \texttt{0x18} & \makecell[l]{\texttt{OVERFLOW}\\\texttt{\_BEHAVIOR}} & чтение &
        Текущее значение параметра \texttt{flush\_logs} \\
        \texttt{0x20} & \texttt{FLUSH} & запись &
        Принудительное чтение доступных сообщений \\
        \texttt{0x28} & \makecell[l]{\texttt{BUF}\\\texttt{SIZE}} & чтение &
        Размер полезной области буфера \\
        \hline
    \end{tabularx}
\end{table}

Для обнаружения устройства операционной системой Stratovirt формирует запись в
таблице ACPI: узел \texttt{LOGD} с идентификатором
\texttt{\_HID = STVRT\_LOG} и ресурсом \texttt{\_CRS}, описывающим диапазон
MMIO-регистров устройства~\cite{acpi_spec}.

\subsection{MPSC-кольцевой буфер}
\label{sec:ch2-ring-buffer}

Общая память организована как кольцевой буфер с несколькими писателями и одним
читателем (MPSC --- Multiple Producer Single Consumer). В
начале буфера располагается заголовок с позициями читателя и писателя, флагом
переполнения и вспомогательной структурой \texttt{SHADOW\_W\_POS}. Эта структура
хранит позицию резервирования для новых сообщений и число писателей, которые
зарезервировали место, но ещё не завершили копирование.

\begin{lstlisting}
typedef union {
    struct {
        uint32_t ReservedPos;
        uint32_t WriterCnt;
    } S;
    uint64_t U;
} SHADOW_W_POS;

typedef struct {
    uint32_t ReadPos;
    uint32_t WritePos;
    uint64_t IsOverflow;
    SHADOW_W_POS ShWPos;
    char Buffer[0];
} LOCKLESS_RING_BUFFER;
\end{lstlisting}

Запись выполняется в три этапа:
\begin{enumerate}
    \item \textbf{reserve} --- писатель атомарно резервирует уникальную область
    буфера;
    \item \textbf{copy} --- сообщение копируется в зарезервированную область;
    \item \textbf{commit} --- писатель сообщает о завершении копирования, после
    чего данные становятся доступными читателю.
\end{enumerate}

На этапе резервирования поток считывает текущее значение
\texttt{SHADOW\_W\_POS}, увеличивает счётчик активных писателей и сдвигает
позицию резервирования на размер сообщения. Обновление выполняется через
атомарную операцию \texttt{compare\_exchange}; при конфликте попытка
повторяется.

\begin{lstlisting}
uint32_t LrbReserve (
    LOCKLESS_RING_BUFFER *Rb,
    const uint32_t MsgSize,
    SHADOW_W_POS *Curr
) {
    for (;;) {
        Curr->U = Rb->ShWPos.U;
        SHADOW_W_POS Next = {
            .S.ReservedPos = Curr->S.ReservedPos + MsgSize,
            .S.WritePos = Curr->S.WriterCnt + 1
        };

        if (atomic_compare_exchange(&Rb->ShWPos, &Curr->U, Next.U)) {
            // SHADOW_W_POS была обновлена в буфере
            // Таким образом, гарантируется, что ни один другой поток
            // не будет копировать своё сообщение в занятую память.
            return Curr->S.ReservedPos;
        }
    }
}
\end{lstlisting}

После копирования сообщения писатель уменьшает счётчик активных писателей. Если
после уменьшения счётчик становится равен нулю, значит все сообщения до текущей
позиции уже скопированы, и глобальная позиция \texttt{WritePos} может быть
продвинута. Читатель видит только область до \texttt{WritePos}, поэтому не
считывает частично скопированные сообщения.

\begin{lstlisting}
void LrbCommit (
    LOCKLESS_RING_BUFFER *Rb,
    SHADOW_W_POS *Curr
) {
    for (;;) {
        SHADOW_W_POS NextShadowPos = {
            .S.ReservedPos = Curr->S.ReservedPos,
            .S.WritePos = Curr->S.WriterCnt - 1
        };

        if (NextShadowPos.S.WriterCnt == 0) {
            // все сообщения до текущего закончили копирование
            Rb->WritePos = Curr->S.ReservedPos;
        }
        if (atomic_compare_exchange(&Rb->ShWPos, &Curr->U, Next.U)) {
            return;
        }
    }
}
\end{lstlisting}

Итоговая функция записи объединяет все три стадии. Сначала она проверяет
возможность записи сообщения, затем резервирует место, копирует сообщение и
фиксирует завершение копирования.

\begin{lstlisting}
uint32_t LrbWrite(
    LOCKLESS_RING_BUFFER *Rb,
    const void *Msg,
    uint32_t MsgSize
) {
    SHADOW_W_POS Curr;
    CONST INT32 ReservedPos = LrbReserve(Rb, MsgSize, &Curr);

    // Копирование сообщения в кольцевой буфер.
    UINT32 WriteChunk = mRingBufferSize - ReservedPos;
    if (WriteChunk > Size) {
        WriteChunk = Size;
    }
    CopyMem(Rb->Buffer + ReservedPos, Msg, WriteChunk);
    if (WriteChunk < Size) {
        CopyMem(Rb->Buffer + ReservedPos, Msg + WriteChunk, Size - WriteChunk);
    }

    LrbCommit(Rb, &Curr);
}
\end{lstlisting}

\subsection{Проверка целостности и переполнения}
\label{sec:ch2-integrity}

Перед резервированием места писатель должен убедиться, что в буфере достаточно
свободного пространства. Для этого вычисляется размер занятой области как
разность между позицией резервирования и позицией читателя. Простая проверка
может ложно принять буфер за повреждённый, если между чтением
\texttt{ShWPos} и чтением \texttt{ReadPos} другой писатель завершил запись, а
читатель успел продвинуть позицию чтения.

Рассмотрим это на примере:
\begin{enumerate}
    \item позиция читателя \texttt{ReadPos = 100};
    \item позиция писателя \texttt{WritePos = 200};
    \item количество активных писателей \texttt{ShWPos.WriterCnt = 1}.
\end{enumerate}

При описанном ниже планировании потоков возникает указанная проблема.

\begin{table}[H]
    \centering
    \adjustbox{max width = \textwidth}{
    \begin{tabular}{|c|c|c|}
        \hline
            Поток-писатель №1 &
            Поток-писатель №2 &
            Поток-читатель \\
        \hline
            \makecell[c]{
                    (1) Операция резервирования: \\
                    Curr->U = Rb->ShWPos.U \\
                    (=> Curr->S.ReservedPos = 200) \\
                    // Поток не увеличил ShWPos.WriterCnt \\
                    // Другие потоки могут обновить WritePos
                } &
            &
            \\
        \hline
            &
            \makecell[c]{
                    (2) Операция резервирования: \\
                    ShWPos.U = Curr->U \\
                    (=> ShWPos.S.ReservedPos = 300)
                } &
            \\
        \hline
            &
            \makecell[c]{
                    (3) Операция commit: \\
                    WritePos = ShWPos.S.ReservedPos = 300 \\
                    // WritePos смог измениться, т.к. \\
                    // поток №1 не увеличил WriterCnt
                } &
            \\
        \hline
            &
            &
            \makecell[c]{
                Читает: \\
                ReadPos = WritePos = 300
                }
            \\
        \hline
            \makecell[c]{
                Операция резервирования: \\
                uint32\_t ReadPos = Rb->ReadPos = 300 \\
                uint32\_t BusySize = Curr->S.ReservedPos - ReadPos = \\
                = 200 - 300
                } &
            &
            \\
        \hline
    \end{tabular}
    }
    \caption{Пример планирования потоков}
    \label{tab:thread-schedule-example}
\end{table}

Таким образом, переменная \texttt{BusySize} переполняется, и буфер ложно
считается повреждённым.

Поэтому при подозрении на неконсистентное состояние значения перечитываются.
Если \texttt{ShWPos} не изменилось, состояние действительно считается ошибочным.
Если изменилось, проверка повторяется с актуальными значениями.

\begin{lstlisting}
uint32_t LrbReserve (
    LOCKLESS_RING_BUFFER *Rb,
    const uint32_t MsgSize,
    SHADOW_W_POS *Curr
) {
    for (;;) {
        Curr->U = Rb->ShWPos.U;
        uint32_t ReadPos = Rb->ReadPos;
        uint32_t BusySize = Curr->S.ReservedPos - ReadPos;

        while (BusySize > mRingBufferSize) {
            CONST UINT64 CurrCpy = Curr->U;
            MemoryFence();
            Curr->U = Rb->ShWPos.U;
            ReadPos = Rb->ReadPos;

            if (Curr->U == CurrCpy) {
                return -1;
            }            
            BusySize = Curr->S.ReservedPos - ReadPos;
        }
        if (BusySize + Size > BUFFER_SIZE) {
            // Сообщение не поместилось.
            // В зависимости от выбранного поведения, сообщение либо
            // пропускается, либо ожидает завершения записи других.
        }

        ...
    }
}
\end{lstlisting}

При недостатке свободного места возможны два режима. В режиме пропуска сообщение
отбрасывается, а флаг переполнения фиксируется в заголовке буфера. В режиме
ожидания писатель инициирует чтение накопленных данных и повторяет попытку
резервирования после освобождения места.

\subsection{Чтение сообщений на стороне хоста}
\label{sec:ch2-reader}

Алгоритм читателя проще алгоритма писателей, так как читатель один. После
получения уведомления через \texttt{KICK} или \texttt{FLUSH} виртуальное
устройство последовательно считывает сообщения от \texttt{ReadPos} до
\texttt{WritePos}. После обработки сообщения позиция \texttt{ReadPos}
продвигается, освобождая место для новых записей.

С точки зрения хоста писатели в буфере не являются доверенными: гостевая ОС
контролирует содержимое памяти, поэтому эмулятор обязан проверять границы
сообщений, размер буфера и состояние указателей. В рамках данной работы
предполагается, что внутри гостевой ОС доступ к интерфейсам записи выдаётся
доверенным компонентам, но это не отменяет проверок на стороне эмулятора.
